\documentclass[10pt,a4paper]{amsart}
\usepackage{amsmath,amssymb,amsthm,graphicx}
% Retained preamble for the Stacks Project chapter snapshot in this directory.
% The chapter's own preamble.tex is a copy of the shared Stacks Project preamble; the block
% below is sliced MECHANICALLY, line for line, from it. Nothing was added, removed or
% redefined; the wrapper class, the xr-hyper/hyperref wiring and the externaldocument list
% are omitted because the assembler supplies the document class.
%
% lines taken from preamble.tex: 3, 5, 9, 11, 12, 156, 157, 158, ... (51 lines total)
%
% package coverage restored from the chapter's own preamble (2026-10-13): xy, xy-2cell

\usepackage[all]{xy}
\xyoption{2cell}
\UseAllTwocells

\newenvironment{reference}{\comment}{\endcomment}
\newenvironment{slogan}{\comment}{\endcomment}
\newenvironment{history}{\comment}{\endcomment}
\theoremstyle{plain}
\newtheorem{theorem}[subsection]{Theorem}
\newtheorem{proposition}[subsection]{Proposition}
\newtheorem{lemma}[subsection]{Lemma}
\theoremstyle{definition}
\newtheorem{definition}[subsection]{Definition}
\newtheorem{example}[subsection]{Example}
\newtheorem{exercise}[subsection]{Exercise}
\newtheorem{situation}[subsection]{Situation}
\theoremstyle{remark}
\newtheorem{remark}[subsection]{Remark}
\newtheorem{remarks}[subsection]{Remarks}
\numberwithin{equation}{subsection}
\def\lim{\mathop{\mathrm{lim}}\nolimits}
\def\colim{\mathop{\mathrm{colim}}\nolimits}
\def\Spec{\mathop{\mathrm{Spec}}}
\def\Hom{\mathop{\mathrm{Hom}}\nolimits}
\def\Ext{\mathop{\mathrm{Ext}}\nolimits}
\def\SheafHom{\mathop{\mathcal{H}\!\mathit{om}}\nolimits}
\def\SheafExt{\mathop{\mathcal{E}\!\mathit{xt}}\nolimits}
\def\Sch{\mathit{Sch}}
\def\Mor{\mathop{\mathrm{Mor}}\nolimits}
\def\Ob{\mathop{\mathrm{Ob}}\nolimits}
\def\Sh{\mathop{\mathit{Sh}}\nolimits}
\def\NL{\mathop{N\!L}\nolimits}
\def\CH{\mathop{\mathrm{CH}}\nolimits}
\def\proetale{{pro\text{-}\acute{e}tale}}
\def\etale{{\acute{e}tale}}
\def\QCoh{\mathit{QCoh}}
\def\Ker{\mathop{\mathrm{Ker}}}
\def\Im{\mathop{\mathrm{Im}}}
\def\Coker{\mathop{\mathrm{Coker}}}
\def\Coim{\mathop{\mathrm{Coim}}}
\DeclareMathSymbol{\boxtimes}{\mathbin}{AMSa}{"02}
\def\QCohstack{\mathcal{QC}\!\mathit{oh}}
\def\Cohstack{\mathcal{C}\!\mathit{oh}}
\def\Spacesstack{\mathcal{S}\!\mathit{paces}}
\def\Quotfunctor{\mathrm{Quot}}
\def\Hilbfunctor{\mathrm{Hilb}}
\def\Curvesstack{\mathcal{C}\!\mathit{urves}}
\def\Polarizedstack{\mathcal{P}\!\mathit{olarized}}
\def\Complexesstack{\mathcal{C}\!\mathit{omplexes}}
\def\Pic{\mathop{\mathrm{Pic}}\nolimits}
\def\Picardstack{\mathcal{P}\!\mathit{ic}}
\def\Picardfunctor{\mathrm{Pic}}
\def\Deformationcategory{\mathcal{D}\!\mathit{ef}}
\usepackage{cleveref}
\crefname{originalproposition}{original proposition}{original propositions}
\Crefname{originalproposition}{Original proposition}{Original propositions}
\makeatletter
\newlabel{definition-hypercovering-variant}{{09VU}{source}{}{sourcepointer-1}{}}
\newlabel{definition-hypercovering-variant@cref}{{[originalproposition][0][]09VU}{[1][0][]source}{}{}{}}
\makeatother
\crefname{originalproposition}{original proposition}{original propositions}
\Crefname{originalproposition}{Original proposition}{Original propositions}
\makeatletter
\newlabel{example-cech}{{01G6}{source}{}{sourcepointer-2}{}}
\newlabel{example-cech@cref}{{[originalproposition][0][]01G6}{[1][0][]source}{}{}{}}
\makeatother
\crefname{originalproposition}{original proposition}{original propositions}
\Crefname{originalproposition}{Original proposition}{Original propositions}
\makeatletter
\newlabel{sets-lemma-coverings-site}{{000X}{source}{}{sourcepointer-3}{}}
\newlabel{sets-lemma-coverings-site@cref}{{[originalproposition][0][]000X}{[1][0][]source}{}{}{}}
\makeatother
\crefname{originalproposition}{original proposition}{original propositions}
\Crefname{originalproposition}{Original proposition}{Original propositions}
\makeatletter
\newlabel{sites-definition-special-cocontinuous-functor}{{03CG}{source}{}{sourcepointer-4}{}}
\newlabel{sites-definition-special-cocontinuous-functor@cref}{{[originalproposition][0][]03CG}{[1][0][]source}{}{}{}}
\makeatother
\crefname{originalproposition}{original proposition}{original propositions}
\Crefname{originalproposition}{Original proposition}{Original propositions}
\makeatletter
\newlabel{theorem-cohomology-hypercoverings}{{01H0}{source}{}{sourcepointer-5}{}}
\newlabel{theorem-cohomology-hypercoverings@cref}{{[originalproposition][0][]01H0}{[1][0][]source}{}{}{}}
\makeatother
\title{Cohomology is the colimit of Cech cohomology over hypercoverings}
\author{The Stacks Project Authors}
\date{stacks-project:hypercovering.tex:local-snapshot-20261012; GNU FDL 1.2 or later}
\begin{document}
\maketitle
\noindent\textit{Editorial note.} Statement span is lines 2402-2421 of sources/stacks-hypercovering/pinned-original-source.tex; proof span is lines 2423-2510, the complete author proof (86 body lines).\par\medskip
\noindent\textit{Editorial note.} Navigation stubs: 5 cited-but-not-included labels are supplied as non-load-bearing pointers with their pinned tags where the snapshot tag index has them.\par\medskip
\noindent\textit{Editorial note.} extra\_packages is empty on purpose: the retained preamble already loads xy with options and a second declaration clashes in the assembler wrapper.\par\medskip
\noindent\textit{Editorial note.} Text is retained verbatim, including the author's own typos and placeholders.\par\medskip
\noindent\textit{Editorial note.} Queue index 2 of queue\_hypercovering.json (pb=4474, stmt 2402-2421, proof 2423-2510).\par\medskip
\noindent\textit{Editorial note.} stacks-hypercovering wave-1 (hc01): fresh adjudication accept (2026-10-14); pool per standing fallback order (hypercovering -> pic -> spaces-more-cohomology) after md01 closed the moduli pool; hypercovering.tex 53 labels/34 with proof judged 26 accept / 8 reject / 14 pb=0 env-dup rows excluded; spans cite-free and reference/slogan/history-free so no patchenv needed. Named omission: Author omits 'the straightforward verifications' of the C' construction and the compatibility of the map into cohomology groups - named mechanical checks; the weak-covering construction and both proofs are in-span.\par\medskip
\noindent\textit{Allowed standard prerequisites (editorial).}\begin{itemize}
\item Explicit prerequisite boundary. The cited inputs are prerequisites; the argument in the retained span is not replaced by any cited result.
\end{itemize}\medskip
\ifdefined\hypertarget\hypertarget{sourcepointer-1}{}\fi
\noindent\textit{Source reference.} Original proposition 09VU: sources/stacks-hypercovering/pinned-original-source.tex, source lines 1285--1298. Cited by the retained statement or proof but not part of the collected spans; supplied only as a navigation pointer. The argument does not depend on its content.\par\medskip
\ifdefined\hypertarget\hypertarget{sourcepointer-2}{}\fi
\noindent\textit{Source reference.} Original proposition 01G6: sources/stacks-hypercovering/pinned-original-source.tex, source lines 587--609. Cited by the retained statement or proof but not part of the collected spans; supplied only as a navigation pointer. The argument does not depend on its content.\par\medskip
\ifdefined\hypertarget\hypertarget{sourcepointer-3}{}\fi
\noindent\textit{Source reference.} Original proposition 000X: sources/stacks-hypercovering/pinned-original-source.tex, source lines 2510--2510. Cited by the retained statement or proof but not part of the collected spans; supplied only as a navigation pointer. The argument does not depend on its content.\par\medskip
\ifdefined\hypertarget\hypertarget{sourcepointer-4}{}\fi
\noindent\textit{Source reference.} Original proposition 03CG: sources/stacks-hypercovering/pinned-original-source.tex, source lines 2510--2510. Cited by the retained statement or proof but not part of the collected spans; supplied only as a navigation pointer. The argument does not depend on its content.\par\medskip
\ifdefined\hypertarget\hypertarget{sourcepointer-5}{}\fi
\noindent\textit{Source reference.} Original proposition 01H0: sources/stacks-hypercovering/pinned-original-source.tex, source lines 2206--2214. Cited by the retained statement or proof but not part of the collected spans; supplied only as a navigation pointer. The argument does not depend on its content.\par\medskip
\begin{proposition}
\label{proposition-cohomology-hypercoverings}
Let $\mathcal{C}$ be a site with fibre products and products of pairs.
Let $\mathcal{F}$ be an abelian sheaf on $\mathcal{C}$.
Let $i \geq 0$. Then
\begin{enumerate}
\item for every $\xi \in H^i(\mathcal{F})$ there exists a hypercovering
$K$ such that $\xi$ is in the image of the canonical map
$\check{H}^i(K, \mathcal{F}) \to H^i(\mathcal{F})$, and
\item if $K, L$ are hypercoverings and $\xi_K \in \check{H}^i(K, \mathcal{F})$,
$\xi_L \in \check{H}^i(L, \mathcal{F})$ are elements mapping
to the same element of $H^i(\mathcal{F})$, then there exists
a hypercovering $M$ and morphisms $M \to K$ and $M \to L$ such
that $\xi_K$ and $\xi_L$ map to the same element of
$\check{H}^i(M, \mathcal{F})$.
\end{enumerate}
In other words, modulo set theoretical issues, the cohomology
groups of $\mathcal{F}$ on $\mathcal{C}$ are the colimit of
the {\v C}ech cohomology groups of $\mathcal{F}$ over all hypercoverings.
\end{proposition}

\begin{proof}
This result is a trivial consequence of
Theorem \ref{theorem-cohomology-hypercoverings}.
Namely, we can artificially replace $\mathcal{C}$ with a slightly
bigger site $\mathcal{C}'$ such that
(I) $\mathcal{C}'$ has a final object $X$ and (II)
hypercoverings in $\mathcal{C}$ are more or less the
same thing as hypercoverings of $X$ in $\mathcal{C}'$.
But due to the nature of things, there is quite a bit of
bookkeeping to do.

\medskip\noindent
Let us call a family of morphisms $\{U_i \to U\}$ in $\mathcal{C}$
with fixed target a {\it weak covering} if the sheafification of the
map $\coprod_{i \in I} h_{U_i} \to h_U$ becomes surjective.
We construct a new site $\mathcal{C}'$ as follows
\begin{enumerate}
\item as a category set $\Ob(\mathcal{C}') = \Ob(\mathcal{C}) \amalg \{X\}$
and add a unique morphism to $X$ from every object of $\mathcal{C}'$,
\item $\mathcal{C}'$ has fibre products as fibre products and products
of pairs exist in $\mathcal{C}$,
\item coverings of $\mathcal{C}'$ are weak coverings of $\mathcal{C}$
together with those $\{U_i \to X\}_{i \in I}$ such that either $U_i = X$
for some $i$, or $U_i \not = X$ for all $i$ and the map
$\coprod h_{U_i} \to *$ of presheaves on $\mathcal{C}$ becomes
surjective after sheafification on $\mathcal{C}$,
\item we apply Sets, Lemma \ref{sets-lemma-coverings-site}
to restrict the coverings to obtain our site $\mathcal{C}'$.
\end{enumerate}
Then $\Sh(\mathcal{C}') = \Sh(\mathcal{C})$ because the inclusion
functor $\mathcal{C} \to \mathcal{C}'$ is a special cocontinuous functor
(see Sites, Definition \ref{sites-definition-special-cocontinuous-functor}).
We omit the straightforward verifications.

\medskip\noindent
Choose a covering $\{U_i \to X\}$ of $\mathcal{C}'$ such that $U_i$ is an
object of $\mathcal{C}$ for all $i$ (possible because
$\mathcal{C} \to \mathcal{C}'$ is special cocontinuous).
Then $K_0 = \{U_i \to X\}$ is a covering in the
site $\mathcal{C}'$ constructed above. We view $K_0$ as an object of
$\text{SR}(\mathcal{C}', X)$ and we set $K_{init} = \text{cosk}_0(K_0)$.
Then $K_{init}$ is a hypercovering of $X$, see
Example \ref{example-cech}. Note that every $K_{init, n}$ has the shape
$\{W_j \to X\}$ with $W_j \in \Ob(\mathcal{C})$.

\medskip\noindent
Proof of (1). Choose $\xi \in H^i(\mathcal{F}) = H^i(X, \mathcal{F}')$
where $\mathcal{F}'$ is the abelian sheaf on $\mathcal{C}'$ corresponding
to $\mathcal{F}$ on $\mathcal{C}$. By
Theorem \ref{theorem-cohomology-hypercoverings}
there exists a morphism of hypercoverings $K' \to K_{init}$
of $X$ in $\mathcal{C}'$ such that $\xi$ comes from an element
of $\check{H}^i(K', \mathcal{F})$.
Write $K'_n = \{U_{n, j} \to X\}$. Now since $K'_n$ maps to
$K_{init, n}$ we see that $U_{n, j}$ is an object of $\mathcal{C}$.
Hence we can define a simplicial object $K$ of $\text{SR}(\mathcal{C})$
by setting $K_n = \{U_{n, j}\}$. Since coverings in
$\mathcal{C}'$ consisting of families of morphisms of $\mathcal{C}$
are weak coverings, we see that $K$ is a hypercovering in the sense
of Definition \ref{definition-hypercovering-variant}.
Finally, since $\mathcal{F}'$ is the unique sheaf on $\mathcal{C}'$
whose restriction to $\mathcal{C}$ is equal to $\mathcal{F}$
we see that the {\v C}ech complexes $s(\mathcal{F}(K))$
and $s(\mathcal{F}'(K'))$ are identical and (1) follows.
(Compatibility with map into cohomology groups omitted.)

\medskip\noindent
Proof of (2). Let $K$ and $L$ be hypercoverings in $\mathcal{C}$.
Let $K'$ and $L'$ be the simplicial objects of $\text{SR}(\mathcal{C}', X)$
gotten from $K$ and $L$ by the functor
$\text{SR}(\mathcal{C}) \to \text{SR}(\mathcal{C}', X)$,
$\{U_i\} \mapsto \{U_i \to X\}$. As before we have equality of
{\v C}ech complexes and hence we obtain $\xi_{K'}$ and
$\xi_{L'}$ mapping to the same cohomology class of $\mathcal{F}'$
over $\mathcal{C}'$. After possibly enlarging our choice
of coverings in $\mathcal{C}'$ (due to a set theoretical issue)
we may assume that $K'$ and $L'$ are hypercoverings of $X$ in
$\mathcal{C}'$; this is true by our definition of hypercoverings in
Definition \ref{definition-hypercovering-variant} and
the fact that weak coverings in $\mathcal{C}$ give coverings in
$\mathcal{C}'$. By
Theorem \ref{theorem-cohomology-hypercoverings}
there exists a hypercovering $M'$ of $X$ in $\mathcal{C}'$
and morphisms $M' \to K'$, $M' \to L'$, and $M' \to K_{init}$
such that $\xi_{K'}$ and $\xi_{L'}$ restrict to the same element of
$\check{H}^i(M', \mathcal{F})$. Unwinding this statement as above
we find that (2) is true.
\end{proof}

\end{document}
